\subsection{Finite tests and comparison of links}

\begin{lemma}[Proved: finite-poset representation of spherical maps]
\label{C:finite-nerve}
Let $E$ be a small category, let $n\geq1$, and let
$f:S^n\to |NE|$ be continuous. There is a finite triangulated
$n$-sphere $K$, an iterated barycentric subdivision $K'$, and a
functor $v:D\to E$, where $D$ is the face poset of $K'$, such that
$Bv$ represents the homotopy class of $f$ under the subdivision
homeomorphism $BD\cong S^n$.
\end{lemma}
\begin{proof}
We use the standard subdivision--extension adjunction
$\operatorname{Sd}\dashv\operatorname{Ex}$. The natural map
$NE\to\operatorname{Ex}^{\infty}NE$ is a weak equivalence and the
target is Kan. Choose the boundary of an ordered $(n+1)$-simplex as the
finite simplicial sphere. In a Kan complex every based $n$th
homotopy class is represented on this boundary by placing an
$n$-simplex with all faces at the basepoint on one face and
placing degenerate basepoint simplices on the other faces.
Thus we can represent the given class on an actual finite
triangulated sphere, whose iterated subdivisions remain
triangulated spheres. Since its set of nondegenerate simplices
is finite, the representative factors through
$\operatorname{Ex}^{r}NE$ for a finite $r$. Its adjoint is a
simplicial map from an iterated subdivision of the chosen sphere to
$NE$. Subdividing once more gives the nerve of a face poset, and
composition with the last-vertex map gives a simplicial map
$ND\to NE$. The nerve functor is fully faithful, so this map is
$Nv$ for a unique functor $v$. The realization of the last-vertex
map is homotopic to the subdivision homeomorphism; the usual
subdivision--extension comparison consequently represents the
original homotopy class. This is an approximation of the map, not
an identification of $E$ with a poset.
\end{proof}

\begin{proposition}[Proved: Lead 6]
\label{C:lead6}
Let $A$ be a small nonempty connected category. For a functor
$u:D\to A$, define the covariant $A$-set
\[
 H_u(a)=\pi_0(u\downarrow a)
\]
and the functor $j_u:D\to\El_AH_u$ by
$j_u(d)=(u(d),[d,\id_{u(d)}])$. Then $A$ is total if and only if
$Bj_u$ is nullhomotopic on every component for every finite poset
$D$ and every such $u$. It suffices to use face posets of finite
triangulated spheres of positive dimension.
\end{proposition}
\begin{proof}
For precision, an object of $u\downarrow a$ is $(d,f:u(d)\to a)$,
and an arrow $(d,f)\to(d',f')$ is $\delta:d\to d'$ with
$f=f'u(\delta)$. Postcomposition defines the $A$-action on $H_u$,
and $j_u(\delta)=u(\delta)$.

Let $H:A\to\Set$ and $v:D\to\El_AH$ be arbitrary. Write
$v(d)=(u(d),h_d)$. Define a natural transformation of $A$-sets
$\lambda:H_u\to H$ by
\[
 \lambda_a[d,f]=H(f)h_d.
\]
The equation $H(u(\delta))h_d=h_{d'}$ makes this independent of
representatives in $\pi_0(u\downarrow a)$; postcomposition proves
naturality. The induced functor
$\El_A\lambda:\El_AH_u\to\El_AH$ satisfies the literal equality
\[
 v=(\El_A\lambda)j_u
\]
on objects and arrows.

If $A$ is total, each component of $\El_AH_u$ has weakly
contractible CW realization and hence contractible realization by
Whitehead's theorem. Therefore each component restriction of $Bj_u$
is nullhomotopic.

Conversely suppose the asserted nullhomotopy condition holds for
sphere face posets. Given $H$, a component $E$ of $\El_AH$, and a
map $S^n\to BE$ with $n\geq1$, Lemma~\ref{C:finite-nerve}
represents its homotopy class by a functor $v:D\to E$. Compose
with $E\hookrightarrow\El_AH\to A$ to obtain $u$. The displayed
factorization and the hypothesis make $Bv$ nullhomotopic. All
homotopy groups of $BE$ therefore vanish. Since $E$ is nonempty
and connected, $BE$ is weakly contractible. This proves totality.
The empty poset adds no condition, and the separate components of
$S^0$ add only the automatic pointwise nullhomotopies.
\end{proof}

\begin{proposition}[Proved: Lead 7]
\label{C:lead7}
Let $A$ be small, $R:A^{\mathrm{op}}\to\Set$ a presheaf, and
$A'$ the category obtained by adjoining one object $Z$ with
$A'(x,Z)=R(x)$, $A'(Z,Z)=\{\id_Z\}$, and no arrows from $Z$ to
old objects. Let $H':A'\to\Set$ and $H=H'|_A$.
The old comma link $L_h$ over $(Z,h)$, for $h\in H'(Z)$, is a
union of components of the two-sided category of elements
$\El(R,H)$.

Suppose a functor $\widetilde u:\widetilde D\to A$ and an
isomorphism
\[
 R\cong\colim_{d\in\widetilde D}A(-,\widetilde u(d))
\]
are specified, and every component of every
$x\downarrow\widetilde u$ is weakly contractible. Then the
canonical functor
\[
 F:\El_{\widetilde D}(\widetilde u^*H)
       \longrightarrow \El(R,H)
\]
is a weak equivalence, as is its restriction over each $h$.
If $n=\dim N\widetilde D<\infty$, then $H_j(L_h;\Z)=0$ for
$j>n$, and
\[
 H_j(\El_AH;\Z)\longrightarrow H_j(\El_{A'}H';\Z)
\]
is injective for $j>n$.
\end{proposition}
\begin{proof}
An object of $\El(R,H)$ is $(x,r,k)$ with $r\in R(x)$ and
$k\in H(x)$. An arrow to $(y,s,l)$ is $f:x\to y$ satisfying
$R(f)s=r$ and $H(f)k=l$. The function
\[
 (x,r,k)\longmapsto H'(r)k\in H'(Z)
\]
is constant along every arrow. Its inverse image of $h$ is
literally $L_h$; it is consequently a union of components.
No injectivity condition on the maps of $H'$ is used.

The specified presheaf presentation identifies
$R(x)=\pi_0(x\downarrow\widetilde u)$. Define
\[
 F(d,k)=(\widetilde u(d),[d,\id],k),\qquad
 F(\delta)=\widetilde u(\delta).
\]
For $e=(x,r,k)$, an object of $e\downarrow F$ is exactly a pair
$(d,f:x\to\widetilde u(d))$ representing $r$: the source
$H$-element in $F$ is forced to be $H(f)k$. The arrow conditions
are exactly those of $x\downarrow\widetilde u$. Thus
$e\downarrow F$ is isomorphic to the component indexed by $r$.
The assumed weak contractibility and Quillen's Theorem A give the
weak equivalence of the specified functor $F$. Restricting to the
union of source components that map to $L_h$ leaves these comma
categories unchanged, so the restriction is a weak equivalence too.

A nonidentity arrow in a category of elements cannot lie over an
identity. Every nondegenerate simplex of
$N\El_{\widetilde D}(\widetilde u^*H)$ therefore projects to a
nondegenerate simplex of $N\widetilde D$. Its dimension is at
most $n$, and the same holds for each union of its components.
This proves the asserted vanishing of link homology.

Every nondegenerate simplex of $N\El_{A'}H'$ is either old or
ends at exactly one new object $(Z,h)$. In the latter case its
preceding chain, including its last arrow into $(Z,h)$, is a
simplex of $NL_h$. Thus the relative chain complex is the direct
sum over $h$ of the augmented link complexes shifted by one.
The long exact sequence contains
\[
 \bigoplus_h H_j(L_h;\Z)\longrightarrow H_j(\El_AH;\Z)
 \longrightarrow H_j(\El_{A'}H';\Z),
\]
which proves injectivity for $j>n$. Infinite sets of new elements
cause no problem because chains are finite sums and homology
commutes with direct sums.
\end{proof}

\begin{remark}[Proved: the limitation of the one-dimensional test]
\label{C:lead7-limitation}
Under Proposition~\ref{C:lead7}'s hypothesis, a necklace diagram
of dimension one or a sum diagram of dimension zero cannot kill
an old $H_2$ class. The comma hypothesis is essential. In a
quasi-filter necklace for $f,g:T\to S$, if $fc=gc$ for an arrow
$c:x\to T$, the copies of $c$ and $fc$ form a circular
$1$-cycle in $x\downarrow\widetilde u$. This comma category
has dimension at most one, so that nonzero cycle cannot be a
boundary. Its component is not weakly contractible. In the
actual-path repair, the explicit identity
$\alpha_{4k}v_1=\alpha_0v_1$ supplies this failure for a subsequent
quasi-filter of $(\alpha_0,\alpha_{4k})$ at $x=V$. Consequently
this proposition does not, by itself, prove persistence through
all stages of the repaired construction.
\end{remark}

\subsection{The interval-type candidate and necessary conditions}

\begin{proposition}[Proved: Lead 8]
\label{C:lead8}
Let $\mathsf{Ty}$ be Grothendieck's category of finite interval
types with endpoint-preserving functors that are nondecreasing
on the numerical vertex order, as in VII manuscript pages
31--32 and PRETOT Definition E.1. Its opposite is the category
of finite words in $+,-$ and weakly increasing sign-preserving
maps of letters. In $\mathsf{Ty}^{\mathrm{op}}$ the two maps
\[
 e=(1,3),\quad c=(1,1):(++)\longrightarrow (+-+)
\]
have no completion to a commutative square. They have no affine
equalizer over any nonzero commutative unital ring. Therefore
$\mathsf{Ty}^{\mathrm{op}}$ is neither total nor exact, and its
colimits fail to be exact over every such ring.
\end{proposition}
\begin{proof}
We first verify the variance and the complete combinatorial
identification. A $\mathsf{Ty}$-map $\tau\to\rho$ has a
nondecreasing vertex function
$\phi:\{0,\ldots,|\tau|\}\to\{0,\ldots,|\rho|\}$ preserving
both endpoints. For every letter $j$ of $\rho$ there is a unique
letter $i$ of $\tau$ such that
$\phi(i-1)<j\leq\phi(i)$. Send $j$ to this $i$.
The resulting map of letters is weakly increasing. It preserves
signs because the nonempty path traversed by a source edge must
have the same sign throughout. Conversely, from a weakly
increasing sign-preserving map $f$ from the letters of $\rho$
to those of $\tau$, define
\[
 \phi(i)=|\{j:f(j)\leq i\}|.
\]
These operations are inverse, also when either word is empty.
The construction commutes with composition: the source edge
traversing a target edge is found by successive pullback of edge
indices. Thus we have an isomorphism of categories on the fixed
word objects, with the stated opposite variance.

Set $p=(++)$ and $q=(+-+)$. For every word $\sigma$ and arrow
$g:q\to\sigma$, monotonicity and the different signs of
adjacent letters imply
$g(1)<g(2)<g(3)$. Thus $ge:p\to\sigma$ is injective,
whereas $hc$ is constant for every $h:q\to\sigma$. Equality
$ge=hc$ is impossible, proving failure of condition I.

Let $R\ne0$ be a commutative unital ring. Define an $R$-linear
functional
\[
 \ell_\sigma:R[\mathsf{Ty}^{\mathrm{op}}(p,\sigma)]\to R
\]
by assigning value $1$ to an injective two-letter map and value
$0$ to a constant two-letter map. These exhaust the possibilities.
For every finite linear combination
$t\in R[\mathsf{Ty}^{\mathrm{op}}(q,\sigma)]$ we have
\[
 \ell_\sigma(te)=\aug(t),\qquad \ell_\sigma(tc)=0.
\]
The equation $te=tc$ therefore implies $\aug(t)=0$, which
excludes augmentation one. Totality implies integral condition
IV by the pair quotient calculation verified in TOTEX A.4, and
exactness over $R$ implies the affine criterion of
\citet[Theorem 3.1 and the remark on p.~296]{IsbellMitchell1976}.
Both conclusions follow. The empty word is initial in
$\mathsf{Ty}^{\mathrm{op}}$, so the category is nonempty and
connected.

We checked the adopted definition against PRETOT E.1 and the
supplied manuscript transcription. At VII ms.~31 the added
restriction is explicitly that the functors are
``croissants sur les objets''. Our independent program
\texttt{code/topology\_ty\_verify.py} compares all $961$ pairs
of sign words of length at most four, finding no discrepancy
between vertex maps and letter maps; checks $2731$ compositions
through length three; and tests $7937$ arrows from $(+-+)$ to
all $1023$ words of length at most nine. None gives a square
completion or overlapping constant/injective supports. These
finite checks supplement the arbitrary-word and arbitrary-ring
proof above.
\end{proof}

\begin{proposition}[Proved: Lead 9]
\label{C:lead9}
Let $A$ be a separator, and let
$a_0,\ldots,a_k:p\to q$ have no integral affine equalizer,
where $k$ is minimal among all such families in $A$. Then
$k\geq2$ and the following hold:
\begin{enumerate}[label=(\alph*)]
\item $A(q,p)=\varnothing$;
\item the cardinalities $|A(p,x)|$ have no finite uniform bound;
\item the full cosieve $C$ on objects receiving an arrow from $p$
is again a separator, containing this family, and $p$ is weakly
initial in $C$;
\item for every arrow $g:q\to x$, the arrows
$ga_0,\ldots,ga_k$ are pairwise distinct.
\end{enumerate}
\end{proposition}
\begin{proof}
Totality supplies I and IV; thus pairs have affine equalizers
and $k\geq2$.

For (a), if $q\to p$ existed,
\citet[Corollary 4.3, p.~297]{IsbellMitchell1976} would equalize
the finite family by an ordinary arrow. We opened the corollary
and its preceding Theorem 4.1 and Corollary 4.2: the second
assertion applies exactly with source $p$ and target $q$, under
$A(q,p)\ne\varnothing$, and requires no weakly terminal object.
The ordinary equalizer would also be affine, a contradiction.

For (b), the representable $A$-set $A(p,-)$ is indecomposable
because its element category $p/A$ has initial object $\id_p$.
A finite bound on its values would allow
\citet[Proposition 4.5, pp.~297--298]{IsbellMitchell1976} to
apply under I and V; IV and V are equivalent by their necklace
criterion. It gives condition II for this representable.
Inductively equalizing pairs then ordinarily equalizes every
finite family in a value, again a contradiction. The source
proof was opened: eventual equality defines a congruence using
I; a quotient value of maximal finite cardinality, together
with the necklace identity, forces all quotient elements equal.
The bound is needed only for this one representable, not
simultaneously for all of them.

For (c), $C$ contains $p,q$ and is a cosieve. Every $C$-set
$H$ extends to an $A$-set by assigning the empty set outside
$C$: there are no arrows from $C$ to its complement, and arrows
entering $C$ receive the unique empty-source function. Its
category of elements is literally $\El_CH$, so totality of $A$
proves totality of $C$. The object $p$ maps to every object of
$C$, hence $C$ is nonempty and connected. Every affine arrow
in $C$ is also one in $A$, so the displayed family remains
unequalizable. Thus $C$ is a separator.

For (d), suppose $ga_i=ga_j$ for $i\ne j$. The set of distinct
arrows among the $ga_l$ has at most $k$ members. By the global
minimality assumption, it has an affine equalizer
$t:x\to y$. The composite $tg:q\to y$ is affine, because
$\aug(tg)=\aug(t)\aug(g)=1$, and it equalizes all $a_l$.
This contradiction proves pairwise distinctness. Global
minimality among families in $A$ is used here; minimality only
among subfamilies of the chosen family would not suffice.
\end{proof}
